\section{The affine level: a per-feature collision frontier}
\label{sec:frontier}
% ============================================================================

Sections~\ref{sec:floor} and~\ref{sec:threshold} bound two criteria on random ensembles. Neither
answers the question an interpretability practitioner actually asks about a \emph{particular}
trained code: for this feature, in this network, up to what sparsity can a linear probe read the
feature's presence off the representation at all? This section gives that quantity exactly, as
the value of one linear programme per feature.

\begin{definition}[Affine separability of a feature]
\label{def:affsep}
Fix $\Phi\in\R^{d\times F}$, a feature $i$ and a sparsity $s$. Write $\mathcal A_i(s)$ for the
set of states $\Phi b$ with $b\in\{0,1\}^F$, $b_i=1$ and $|\operatorname{supp}(b)|\le s$, and
$\mathcal B_i(s)$ for the same with $b_i=0$. Feature $i$ is \emph{affinely separable at sparsity
$s$} if there exist $w\in\R^d$ and $\theta\in\R$ with $w^{\top}x>\theta$ for every
$x\in\mathcal A_i(s)$ and $w^{\top}x<\theta$ for every $x\in\mathcal B_i(s)$.
\end{definition}

The states are Boolean here, and that is not a simplification we could drop. Allowing an active
amplitude to be any positive number makes ``$i$ is active'' an open condition, and an active
state with $b_i\to0$ converges to an inactive one; the two hulls then touch for every code and
every sparsity, so nothing would ever be separable. A positive floor on the active amplitude is
what makes the question well posed, and Corollary~\ref{cor:kappa-alpha} is the graded statement
with that floor made explicit.

An affine functional separates two sets exactly when it separates their convex hulls, so the
question is whether those two hulls meet. The programme below decides it.

\begin{theorem}[The collision frontier]
\label{thm:frontier}
For $z\in\R^{F-1}$, indexed by the features other than $i$, let $z^{+}_j=\max(z_j,0)$ and
$z^{-}_j=\max(-z_j,0)$, so that both parts are non-negative and $z=z^{+}-z^{-}$. Let
$\Phi_{-i}$ be $\Phi$ with column $i$ deleted. Put
\begin{equation}
  \kappa_i\;:=\;\min_{z}\Big\{\max\big(\textstyle\sum_j z^{+}_j,\;1+\sum_j z^{-}_j\big)
  \;:\;\Phi_{-i}z=\phi_i,\;\norm{z}_\infty\le1\Big\}.
  \label{eq:kappa}
\end{equation}
with $\kappa_i:=+\infty$ when the constraints admit no $z$ at all. Then feature $i$ is affinely
separable at every sparsity $s'\le s$ if and only if $\kappa_i>s$. When $\kappa_i$ is finite the
largest sparsity at which feature $i$ is separable is $\lceil\kappa_i\rceil-1$; when
$\kappa_i=+\infty$ the feature is separable at every sparsity, and no collision exists to bound.
\end{theorem}

\begin{proof}
Separability fails exactly when $\operatorname{conv}\mathcal A_i(s)$ meets
$\operatorname{conv}\mathcal B_i(s)$. By Proposition~\ref{prop:knapsack} below, applied to the
coordinates other than $i$ with budget $s-1$ on the active side and $s$ on the inactive side,
those hulls are $\Phi\{e_i+u:u\in[0,1]^{F-1},\;\mathbf 1^{\top}u\le s-1\}$ and
$\Phi\{v:v\in[0,1]^{F-1},\;\mathbf 1^{\top}v\le s\}$ respectively --- the relaxation from
$0/1$ coordinates to the box is exact, because the polytope is integral. So a collision is a
pair $u,v$ in those two boxes with $\Phi(e_i+u)=\Phi v$. We may take $u$ and $v$ to have
disjoint supports, since a common component cancels from both sides and only relaxes the two sum
constraints. Then
$\Phi(e_i+u)=\Phi v$ reads $\Phi_{-i}(v-u)=\phi_i$, so with $z=v-u$ we have $z^{+}=v$ and
$z^{-}=u$. The budget on the inactive state is $\sum_j v_j\le s$, that on the active state is
$1+\sum_j u_j\le s$, and the entries lie in $[0,1]$, giving $\norm{z}_\infty\le1$. The two
budgets are exactly $\sum_jz^{+}_j\le s$ and $1+\sum_jz^{-}_j\le s$, so a collision at some
sparsity at most $s$ exists if and only if the objective in~\eqref{eq:kappa} can be brought to
$s$ or below, that is, if and only if $\kappa_i\le s$.
\end{proof}

Three things improve on the corresponding exactly-$s$ statement, which we record because the
earlier formulation of this work used it. The threshold is $s$ itself rather than
$2\min(s,F-s)-1$; the admissible sets nest in $s$, so monotonicity of the frontier holds by
construction rather than needing a proof and a restriction to $s\le F/2$; and $\kappa_i=7.4$ now
reads directly as ``separable up to $7$, lost at $8$'', which is what a capacity ought to mean.
Computationally~\eqref{eq:kappa} is a single linear programme, stated as
Algorithm~\ref{alg:frontier}: the $\max$ of two linear forms is handled by one auxiliary scalar,
and one solve settles every sparsity at once rather than one feasibility problem per $s$.

\begin{algorithm}[t]
\caption{\textsc{CollisionFrontier} --- the largest affinely separable sparsity, per feature}
\label{alg:frontier}
\begin{algorithmic}[1]
\Require code $\Phi\in\R^{d\times F}$; feature set $\mathcal I\subseteq[F]$; minimum amplitude
  $\alpha\in(0,1]$
\Ensure $\kappa_i(\alpha)$ and the frontier $\lceil\kappa_i(\alpha)\rceil-1$ for each
  $i\in\mathcal I$
\For{$i\in\mathcal I$}
  \State solve, in variables $u,v\in\R^{F-1}_{\ge0}$ and $t\in\R$:
  \Statex \qquad $\min\;t$ \quad s.t.\quad $\Phi_{-i}(u-v)=\phi_i$,\quad
    $\alpha\mathbf 1^{\top}u\le t$,\quad $1+\alpha\mathbf 1^{\top}v\le t$,\quad
    $u,v\le1/\alpha$
  \State $\kappa_i(\alpha)\gets t^{\star}$;\quad
    $s_i^{\max}\gets\lceil\kappa_i(\alpha)\rceil-1$
\EndFor
\State \Return $\{(\kappa_i(\alpha),s_i^{\max})\}_{i\in\mathcal I}$
\end{algorithmic}
\end{algorithm}

\begin{remark}[Why splitting $z$ into two non-negative vectors loses nothing]
\label{rem:split}
Algorithm~\ref{alg:frontier} does not constrain $u$ and $v$ to have disjoint supports, so a
feasible point need not have $u=z^{+}$ and $v=z^{-}$. It is still exact. If
$u_j,v_j>0$ for some $j$, subtracting $\min(u_j,v_j)$ from both leaves $u-v$ unchanged, keeps
feasibility, and strictly decreases both $\mathbf 1^{\top}u$ and $\mathbf 1^{\top}v$, hence does
not increase $t$. Every optimum therefore admits a canonical split, on which the two budget
constraints are exactly those of~\eqref{eq:kappa} and the bounds $u,v\le1/\alpha$ are exactly
$\norm{z}_\infty\le1/\alpha$.
\end{remark}

The cost is one linear programme in $2(F-1)+1$ variables with $d$ equality constraints per
feature, which is affordable for every feature at $d=50$ and not at $d=400$; how we choose a
subset there, and why the resulting minimum is an upper bound, is stated in
Section~\ref{sec:limitations}.

\subsection{What the state distribution contributes}
\label{sec:frontier-amplitudes}

Definition~\ref{def:affsep} allowed graded activations in $[0,1]$. It is natural to ask what
happens under a genuine state law, and the answer is unusually clean: almost nothing survives.

\begin{proposition}[Amplitudes wash out of the hulls]
\label{prop:knapsack}
For $\alpha\in(0,1]$ and $k\in\mathbb N$ let
$\mathcal U_k^{\alpha}=\{u\in\R_{\ge0}^{F}:u_j\in\{0\}\cup[\alpha,1],\;
|\operatorname{supp}(u)|\le k\}$. Then
$\operatorname{conv}\mathcal U_k^{\alpha}=\{u\in[0,1]^F:\mathbf 1^{\top}u\le k\}$ for every
$\alpha$.
\end{proposition}

\begin{proof}
The right-hand side is the unit-weight knapsack polytope, which is integral: its vertices are the
$0/1$ vectors with at most $k$ ones. Each such vertex lies in $\mathcal U_k^{\alpha}$ for every
$\alpha\le1$, so the right-hand side is contained in the hull; the reverse inclusion is immediate
since $\mathcal U_k^{\alpha}$ satisfies both constraints.
\end{proof}

\begin{corollary}[The frontier under a minimum amplitude]
\label{cor:kappa-alpha}
If active amplitudes lie in $[\alpha,1]$ then feature $i$ is separable at every sparsity
$s'\le s$ iff $\kappa_i(\alpha)>s$, where
\[
  \kappa_i(\alpha)=\min_z\Big\{\max\big(\alpha\textstyle\sum_jz^{+}_j,\;
  1+\alpha\sum_jz^{-}_j\big)\;:\;\Phi_{-i}z=\phi_i,\;\norm{z}_\infty\le1/\alpha\Big\},
\]
and $\kappa_i(\alpha)$ is non-decreasing in $\alpha$. As $\alpha\to0$ an active state converges
to an inactive one, the hulls touch, and no uniform margin exists.
\end{corollary}

\begin{remark}[The affine level is a support-level worst case, not a weighted average]
\label{rem:notdistributional}
It would be convenient to call Corollary~\ref{cor:kappa-alpha} a distribution-aware frontier. It
is not, and the distinction matters. By Proposition~\ref{prop:knapsack} two state laws with the
same support family and the same minimum amplitude give the \emph{same} $\kappa_i(\alpha)$,
however differently they weight the states in between. The affine level therefore reads exactly
one scalar off the state law --- the smallest detectable amplitude --- whereas by
Proposition~\ref{prop:secondmoment} the analog level reads the full second moment $C$. That
asymmetry is a property of the two criteria, not an artefact of our formulation: separability is
a statement about a worst case over a support family, and a worst case cannot be reweighted. A
genuinely probabilistic version, guaranteeing separability for all but a $\delta$ fraction of
states, is possible but needs machinery we do not open here.
\end{remark}

\subsection{Margin, certificate, and a dimensional ceiling}
\label{sec:frontier-margin}

The verdict of Theorem~\ref{thm:frontier} is binary. Two further quantities say how comfortably
it holds and how it fails, and one says that no code can do well beyond a certain sparsity. The
first two are stated for the exactly-$s$ parameterisation in which they were derived and
verified, where the admissible set carries the affine constraint $\mathbf 1^{\top}z=1$ and an
$\ell_1$ budget; we flag that rather than restate them without proof.

\begin{proposition}[The robust margin is half the hull distance]
\label{prop:margin}
Let $\delta_i(s)=\operatorname{dist}(\mathcal A_i(s),\mathcal B_i(s))$ be the Euclidean distance
between the two hulls. The maximum-margin affine decoder for feature $i$ has
$\norm{w^{\star}}=1/\delta_i(s)$ and margin $\gamma_i(s)=1/(2\norm{w^{\star}})=\delta_i(s)/2$.
Consequently, with the optimal unit-norm $w$ and its centred bias, every perturbed input
$x+\eta$ with $\norm{\eta}_2<\gamma_i(s)$ still receives the correct label on every admissible
support.
\end{proposition}

\begin{proof}
The margin programme is $\min\norm{w}$ subject to $\min_{u\in D}w^{\top}u\ge1$ with
$D=\mathcal A_i(s)-\mathcal B_i(s)$. Support-function duality gives
$\max_{\norm{w}\le1}\min_{u\in D}w^{\top}u=\operatorname{dist}(0,D)=\delta_i(s)$, whence
$\norm{w^{\star}}=1/\delta_i(s)$. For the perturbation claim, the score moves by at most
$\norm{w}\norm{\eta}$.
\end{proof}

The margin is therefore a certified noise tolerance in the units of the representation, not
merely a scale, and $\delta_i(s)=0$ holds exactly when the feature is not separable --- the
verdict and the margin are two readings of one object.

\begin{proposition}[Failure comes with a finite, independently checkable certificate]
\label{prop:farkas}
If feature $i$ is not separable at sparsity $s$, there are states
$x^{+}_1,\dots,x^{+}_m\in\mathcal A_i(s)$, states $x^{-}_1,\dots,x^{-}_n\in\mathcal B_i(s)$ and
convex weights $\lambda,\mu$ with
\begin{equation}
  \sum_{k}\lambda_kx^{+}_k\;=\;\sum_{l}\mu_lx^{-}_l .
  \label{eq:farkas}
\end{equation}
No affine rule can label all of them correctly, since affinity would force the same point to lie
strictly on both sides of $\theta$. By Radon's theorem the obstruction can be compressed to at
most $d+2$ states.
\end{proposition}

\begin{remark}[What the certificate is, and what it is not]
\label{rem:cert}
Equation~\eqref{eq:farkas} is a convex combination of active states equal to a convex combination
of inactive ones. It is \emph{not} a single pair of colliding supports, and describing it as one
would be wrong: the obstruction is generically spread over several states on each side. This
matters practically, because the certificate is what makes the negative verdict auditable. ``The
solver reported infeasible'' is not evidence; a list of supports from which any third party can
rebuild the states from $\Phi$ and recompute~\eqref{eq:farkas} is. Our implementation stores the
supports rather than the cut vectors for exactly that reason, and re-verifies them in a process
that never sees the solver.
\end{remark}

\begin{proposition}[A dimensional ceiling on the affine level]
\label{prop:circuit}
Call a set $C$ of columns an \emph{affine circuit} if it is minimally affinely dependent:
$\sum_{j\in C}l_j\phi_j=0$ with $\sum_{j\in C}l_j=0$ and every $l_j\neq0$. Let $q$ be the
smallest size of an affine circuit of $\Phi$. Since any $d+2$ points in $\R^d$ are affinely
dependent, $q\le d+2$ whenever $F>d+1$. Moreover the exactly-$s$ collision radius satisfies
$\rho_{\min}\le q-1\le d+1$, and hence no code supports featurewise affine support recovery
beyond sparsity $\lceil d/2\rceil$.
\end{proposition}

\begin{proof}
Given a circuit $C$, pick $i\in C$ with $|l_i|$ largest and set $z_j=-l_j/l_i$ for
$j\in C\setminus\{i\}$ and $z_j=0$ otherwise. Then $\Phi_{-i}z=\phi_i$; the affine constraint
$\mathbf 1^{\top}z=1$ holds automatically, because $\sum_{j\in C}l_j=0$ forces
$\sum_{j\neq i}l_j=-l_i$; $\norm{z}_\infty\le1$ by the choice of $i$; and
$\norm{z}_1\le q-1$. Hence $\rho_i\le q-1$. The sparsity bound follows from the threshold
$2\min(s,F-s)-1$ of the exactly-$s$ rule.
\end{proof}

This is the affine counterpart of Theorem~\ref{thm:floor}: one bounds error from below, the other
bounds sparsity from above, and both follow from the ambient dimension alone.

% ============================================================================
